\documentclass[11pt,a4paper]{article}
\usepackage[T1]{fontenc}
\usepackage{isabelle,isabellesym}
\usepackage{amsmath,amssymb,url}

% This should be the last package used.
\usepackage{pdfsetup}

\urlstyle{rm}
\isabellestyle{it}

\begin{document}

\title{Poisson Summation and the Jacobi Theta Transformation}
\author{Lawrence C. Paulson}
\date{}
\maketitle

\begin{abstract}
This entry formalises the one-dimensional Poisson summation formula for
continuous complex-valued functions. The main theorem assumes a summable
majorant for the translates on one unit interval and absolute summability
of the integer Fourier samples. It also gives a criterion in terms of power
decay and a version for arbitrary positive lattice spacing. Examples
compute the Fourier samples of the tent function and derive the Jacobi
theta transformation from the Fourier transform of a Gaussian.
\end{abstract}

\tableofcontents

\section{Overview and sources}

For the Fourier transform
\[
  \widehat f(\xi)=\int_{\mathbb R} f(t)e^{-2\pi i\xi t}\,dt,
\]
the principal result identifies the periodisation
$\sum_{n\in\mathbb Z} f(x+n)$ with its Fourier series. The lattice-spacing
theorem replaces the unit lattice by $h\mathbb Z$ for $h>0$ and has Fourier
frequencies $k/h$. The Gaussian example yields
\[
  \sqrt t\,\sum_{n\in\mathbb Z}e^{-\pi t n^2}
  =\sum_{k\in\mathbb Z}e^{-\pi k^2/t}
  \qquad (t>0).
\]

Poisson summation and its Gaussian application are standard results in
Fourier analysis~\cite{stein_shakarchi_fourier}. The proof architecture
follows the coefficient calculation for a periodisation and Fourier
reconstruction used in the Mathlib formalisation by
Loeffler~\cite{mathlib_poisson}. Reconstruction here uses the Fej\'er
convergence theorem from the AFP entry \emph{Fourier Series}~\cite{afp_fourier}.
The Gaussian transform is obtained from the standard normal
characteristic-function theorem in Isabelle/HOL's
\texttt{HOL-Probability.Characteristic\_Functions}.

\input{session}

\bibliographystyle{abbrv}
\bibliography{root}

\end{document}
